\section{Geração de princípios e combinações exploratórias}
A função global Proteger louça se desdobra em sete funções parciais, 21 elementares e sete complementares. Na matriz morfológica, cada linha representa uma função de menor nível e cada coluna apresenta um mecanismo possível. E04.1 identifica a primeira opção para absorver impacto. Para cada função, C1, C2 e C3 registram a opção usada em cada combinação.

Para combinar os mecanismos, é preciso verificar os contatos com a cerâmica, o caminho das cargas, a vedação durante a inspeção e as interfaces de montagem e manutenção. Uma mesma peça pode atender a várias funções. A análise de interfaces determina como os princípios trabalham juntos.

A matriz morfológica organiza alternativas por função e pode registrar princípios por descrição verbal ou representação gráfica \citep[pp.~266, 269]{back2008}. O Anexo~\ref{anexo:matriz} preserva as 112 opções ilustradas e as escolhas C1/C2/C3 da base da equipe \citep{baseequipe2026}. As imagens são referências de mecanismos e aplicações análogas, não comprovações do desempenho do dispositivo.

\subsection{Combinações exploratórias}
Cada combinação reúne 28 escolhas. A matriz no anexo registra as três escolhas ao lado de cada função, da primeira elementar à última complementar.

\begin{singlespace}
\small
\begin{longtable}{@{}L{0.06\linewidth} L{0.46\linewidth} L{\dimexpr0.48\linewidth-4\tabcolsep\relax}@{}}
\caption{Síntese das combinações exploratórias}\label{tab:combinacoes}\\
\toprule
\textbf{ID} & \textbf{Princípios centrais} & \textbf{Interfaces a verificar} \\
\midrule\endfirsthead
\multicolumn{3}{l}{\small Quadro \thetable{} — continuação}\\
\toprule\textbf{ID} & \textbf{Princípios centrais} & \textbf{Interfaces a verificar} \\
\midrule\endhead
\midrule\multicolumn{3}{r}{\small Continua na próxima página.}\\\endfoot
\bottomrule\endlastfoot
C1 & Berço com amortecimento celular, cintas limitadas, guias ortogonais e carrinho externo. Envoltório com visores, bandeja retida e colunas periféricas. & Conferir tensão e área de contato, partida no carrinho, transparência após desgaste e recolhimento das paredes entre as colunas. \\ \addlinespace[4pt]
C2 & Quadro telescópico, sapatas flutuantes, elastômeros e rolamento removível. Painéis transparentes, cuba coletora e conexões indexadas. & Conferir limpeza dos telescópios, massa, volume carregado, desvio da carga para fora da cerâmica e alinhamento após reposição. \\ \addlinespace[4pt]
C3 & Painéis articulados, apoios conformáveis e mobilidade integrada. Bolsa coletora, visores distribuídos e juntas em labirinto. & Conferir freios e apoios no piso, ângulo de tombamento, transferência assistida e separação das camadas ao fim de vida. \\ \addlinespace[4pt]
\end{longtable}
\fonte{Elaboração da Equipe 2; conteúdo das três combinações preservado da base de 4 de novembro de 2026.}
\end{singlespace}
\subsection{Verificações para a próxima seleção}
O apoio precisa imobilizar a louça e, ao mesmo tempo, manter curso disponível para absorver impacto. A travessa e as colunas devem conduzir a compressão fora da cerâmica. A barreira precisa permitir inspeção com o conjunto fechado e reter fragmentos mesmo depois de uma troca de módulo.

A comparação entre as combinações depende de verificar, nas mesmas condições de uso, as metas de impacto, deslocamento, inspeção, contenção e movimentação. A seleção da concepção física ocorrerá depois do estudo de módulos e interfaces. As verificações propostas neste relatório são atividades futuras, não resultados de ensaios já realizados.

\subsection{Ressalva econômica preliminar}
O custo por ciclo do dispositivo compreende aquisição rateada pela vida efetiva, recondicionamento, retorno e consumíveis. Para comparar M1 com a situação atual, é necessário acrescentar o preparo e a transferência no distribuidor, os custos atribuíveis à retirada e à destinação da embalagem e a avaria residual de cada cenário, mantendo a mesma fronteira operacional.

O valor de R\$~15 por ciclo em RP23 permanece como meta provisória, não como comprovação de viabilidade. O custo total da alternativa deve ser comparado ao custo total da situação atual; a economia de avarias e eventuais economias de embalagem não podem ser contadas em duplicidade. Nesta etapa, cabe uma ressalva de coerência e acompanhamento econômico, não um estudo financeiro completo \citep[pp.~339--340]{back2008}.
